\chapter{CONCLUSION}

This chapter summarizes the thesis, states the remaining limitations, and outlines future work. Because the topic is student mental-health support, the conclusion avoids clinical claims and treats the system as a research prototype.

\section{Summary of Contributions}

This thesis proposed a therapist-led virtual group therapy prototype for student mental-health support. The system combines DASS-21-based onboarding, route-aware intervention, blackboard coordination, Yalom-inspired peer voices, therapist planning, safety validation, and session memory. Instead of generating a direct chatbot reply, the system first interprets the user's situation, selects a therapeutic route, controls peer participation, and then produces a final therapist response.

The main contribution is a structured design for controlled group-style support. Nam and Linh are not independent therapists and do not control the session. Their role is limited to short peer contributions that may support universality, catharsis, hope, or interpersonal learning. The therapist core remains responsible for the session flow, safety boundaries, and final response. This makes the system closer to a facilitated virtual group than to a free multi-agent chat.

The second contribution is the VSM evaluation framework. VSM-All provides synthetic multi-turn student-support cases for testing route fidelity, conversation continuity, peer timing, safety boundaries, and overall response quality. The full benchmark shows that the proposed system performs better than the evaluated baselines, especially in safety, group dynamics, continuity, and structured therapeutic support.

\section{Limitations}

The results should be interpreted carefully. Although the full VSM-All benchmark was conducted, the evidence remains prototype-level because the cases are synthetic and the evaluator is LLM-based. The benchmark is useful for structured comparison and regression testing, but it cannot replace clinician review or real-user evaluation.

The system also has clinical limitations. It has not been tested in a clinical trial, has not been validated by licensed therapists on real student users, and should not be treated as a treatment provider. It may still produce generic, repetitive, poorly timed, or incomplete responses, especially in CBT cases that require careful Socratic questioning and cognitive restructuring.

Finally, the architecture is more complex than a single-prompt chatbot. The blackboard, peer drafts, therapist planner, validators, and memory updater make the system more controllable, but they can also increase latency and token usage. A real deployment would need stronger optimization, monitoring, and human oversight.

\section{Ethical Boundaries and Future Work}

Mental-health support systems require strict boundaries. The prototype should not diagnose users, give medication advice, promise recovery, replace professional care, or encourage emotional dependence on the system. Peer voices also require caution because they can make the interaction feel socially real. For this reason, Nam and Linh must remain limited, transparent, and therapist-governed.

Future work should focus on three directions. First, the benchmark should be validated with human and clinician review so that LLM-judge scores can be calibrated against expert judgment. Second, the CBT route should be strengthened with better stage evidence, better guided-discovery questions, and stricter checks against premature advice. Third, the system should be optimized for lower latency and safer long-term memory before any real user-facing deployment is considered.

\section{Conclusion}

This thesis does not claim that the system provides clinical treatment. Its contribution is a controllable architecture and evaluation workflow for studying virtual group-style mental-health support. The results suggest that therapist-led coordination, bounded peer participation, and explicit safety checking can make LLM-based support more structured than a generic chatbot response. With clinician audit, real-user studies, and further safety work, the prototype can serve as a useful research platform for student-support dialogue.
